\documentclass[12pt]{article}

% ── Packages ──────────────────────────────────────────────────────────────
\usepackage[margin=1in]{geometry}
\usepackage{amsmath, amssymb}
\usepackage{graphicx}
\usepackage{booktabs}
\usepackage{natbib}
\usepackage{hyperref}
\usepackage{xcolor}
\usepackage{soul}          % \hl{} for TODO highlights
\usepackage{enumitem}
\usepackage{array}
\usepackage{caption}
\usepackage{subcaption}
\usepackage{multirow}
\usepackage{lineno}        % line numbers for review
\linenumbers

% ── TODO command ──────────────────────────────────────────────────────────
\newcommand{\todo}[1]{\hl{\textbf{[TODO: #1]}}}
\newcommand{\note}[1]{{\color{blue}\textit{[Note: #1]}}}

% ── Title block ───────────────────────────────────────────────────────────
% Working title -- submission venue not yet decided. Renamed from
% draft_informs_2026.tex; retitled to foreground the patient-harm/adverse-event
% story (the main narrative) over the extraction pipeline (a component within it).
\title{\textbf{From Regulatory Text to Patient Harm:\\
  LLM Extraction of FDA Form 483 Signals for Adverse Event Prediction}}

\author{
  Amirreza Sahebi-Fakhrabad\thanks{Corresponding author.
    Poole College of Management, North Carolina State University.
    Email: \texttt{asahebi@ncsu.edu}} \and
  Shailesh Divey\thanks{Fisher College of Business, The Ohio State University.} \and
  John V.\ Gray\footnotemark[2] \and
  Robert Handfield\footnotemark[1] \and
  Amir Sadeghi\footnotemark[1]
}

\date{\today \\ \vspace{4pt}
  \textit{Working draft. Submission venue not yet finalized.
  Please do not cite or circulate without permission.} \\
  \vspace{4pt}
  \note{Results, Discussion, and Conclusion are pending: the extraction
  pipeline was substantially revised after an external expert review
  (Section~\ref{sec:methods}), and a blind human-labeler validation round is
  in progress (n=50 pilot underway; a larger stratified round planned).
  The numbers in those sections below are from the prior pipeline version
  and are placeholders showing the intended shape of the analysis, not
  citable results. Introduction and Methodology reflect the current pipeline.}}

% ══════════════════════════════════════════════════════════════════════════
\begin{document}
\maketitle

% ── Abstract ──────────────────────────────────────────────────────────────
\begin{abstract}
Manufacturing quality failures at generic drug facilities put patients at
risk, but FDA generates a lot of unstructured regulatory text on these
facilities that goes almost entirely unused as a predictive resource:
inspection observations, warning letters, recall narratives, import refusal
records.
This paper builds and validates an LLM pipeline that turns FDA Form 483
inspection observations into structured, patient-harm-relevant signals, for
129 pharmaceutical manufacturing facilities across 14 high-priority generic
Active Pharmaceutical Ingredients (APIs).
We use NLP on the structured regulatory data and an LLM pipeline (Claude,
with a strict JSON schema, calibrated with an external pharmaceutical-quality
expert) to process the Form 483 observation text.
This gives interpretable features for violation type, graded severity,
failure scope, escalation pattern, and firm behavior.
Our main claim is what we call the silent-problem pattern.
FDA's three-class inspection outcome (OAI/VAI/NAI) hides more than it
reveals, particularly inside the VAI class where most facilities in our data
land, because the internal review that produces the label (the Establishment
Inspection Report and the facility's response to it) is not visible to
researchers.
483 text is closer to a direct record of what the inspector actually found.
\todo{Re-check this claim once we have current-pipeline numbers -- may need
softening or strengthening depending on what the VAI-heterogeneity analysis
actually shows.}
We test whether that text signal predicts facility-level adverse events, and
whether AE trajectories diverge after an OAI outcome (mandatory corrective
action) versus a VAI or NAI outcome (no requirement).
That is our working hypothesis, not yet an established finding.
Preliminary results, from an earlier version of the pipeline that predates
the calibration described in this draft, are encouraging: specific
regulatory citation patterns strongly predict the most severe inspection
outcomes (Pearson $r = 0.51$), and LLM severity grading separates the most
severe inspection outcomes (Critical+Major share 20 percentage points higher
in Official-Action-Indicated inspections than Voluntary-Action-Indicated).
\todo{Replace with current-pipeline numbers once the Stage 3 validation and
full-corpus run are complete.}
\end{abstract}

\textbf{Keywords:} FDA Form 483, manufacturing quality, adverse events,
large language model, NLP, pharmaceutical supply chain

\newpage

% ══════════════════════════════════════════════════════════════════════════
\section{Introduction}
\label{sec:intro}

Generic drugs account for roughly 90\% of U.S. prescriptions filled
\citep{fda2023generic}, yet generic drug shortages are an enduring public
health problem that disrupts patient care, inflates costs, and strains
hospital operations \citep{fox2014shortage}.
Manufacturing quality failure is a leading proximate cause of these
shortages, not just an incidental one \citep{NaumovJOOM2025,
OorschotDecSci2022}: shortages arise from an interplay of demand
concentration, supply chain fragility, and latent quality failures that may
not surface until facilities fail FDA inspections or recall products.

FDA inspects a given manufacturing facility only from time to time.
Between visits, nobody outside the facility has real-time insight into what
is happening on the floor.
Quality can drift for a long time before anyone catches it, and by the time
it surfaces as a recall or an enforcement action, patients have often
already been affected.

Here is how an inspection actually ends.
On the last day, the inspector hands the facility a Form 483: a numbered
list of every Good Manufacturing Practice (GMP) deficiency observed, written
in the inspector's own words. It covers contamination controls, laboratory
records, data integrity, quality-system gaps, whatever was actually found at
that facility on that day.
After returning to the office, the same inspector writes a much larger
internal document, the Establishment Inspection Report (EIR): the full
narrative behind each observation, supporting exhibits, batch records,
management discussion, and the facility's written response to the 483.
FDA uses the EIR and that response, not the 483 text alone, to decide the
final classification. No Action Indicated (NAI), Voluntary Action Indicated
(VAI), and Official Action Indicated (OAI) are the three possible outcomes,
and the classification goes to the facility separately, later.

This creates a specific data problem for anyone trying to study inspection
quality.
Researchers can obtain the 483 text, in our case through Redica's licensed
FOIA-based database (Section~\ref{sec:data}).
The EIR, which is richer and contains the facility's actual response, is not
systematically available to researchers.
The final classification is the only fully public signal.
That means we do not know how FDA gets from the 483 text to the
classification. The EIR and the internal review behind it are invisible to
us.
The relationship between what the 483 says and what label a facility
ultimately receives is itself an open question, and there is documented
reason to think it is not a clean function of severity alone: investigative
reporting has described cases where a facility's classification was
downgraded from OAI to VAI for reasons outside the inspection findings,
including concern over drug shortages \citep{OAIVAIDowngradeReporting,
GrayEtAlUnpublished}.

Relying only on the three-class outcome label while ignoring the inspector's
own narrative text throws away most of what the inspection actually
produced.
The gap between what the label says and what the text says is what this
paper studies.

Our working hypothesis, the silent-problem pattern, applies only to
facilities with a real, developing quality problem.
At those facilities, adverse events should start rising before the
inspection catches the problem, plausibly in part because FDA's own
site-selection process weighs adverse-event signals when deciding whom to
inspect.
The inspection follows, and the underlying problem gets written down as 483
text.
FDA then assigns an outcome.
If that facility is classified OAI, corrective action is mandatory and we
would expect adverse events to fall.
If it is classified VAI or NAI, nothing is required, and we would expect
elevated adverse events to persist, even though VAI and NAI both look, on
paper, like the facility is basically fine.
Text severity should let us tell apart the two kinds of VAI facility: the
one with a genuinely serious, unresolved problem, and the one that is
actually fine, a distinction the label alone cannot make.
\note{Check: is the ae\_validation-folder analysis (VAI-only AUC, Hi-sig
vs.\ Lo-sig VAI adverse-event trajectories) a validated, current result, or
still built on the stale pre-Redica pipeline? If current, upgrade this from
hypothesis to reported finding and move the actual numbers into the
Abstract and Results.}

Prior work shows FDA inspection outcomes predict later drug shortages
\citep{Stomberg2018, Wang2025MSOM}, and that machine learning can forecast
shortages from pharmacy demand data \citep{Liu2021, Pall2023}.
But nobody has used the narrative text inspectors write in 483 observations
as a predictive input, and no study has linked inspection-text signals
directly to facility-level patient harm through adverse event reports
\citep{Pazhayattil2020}.

The LLM pipeline is the tool that makes the 483 text measurable and
comparable across facilities. It is not the contribution by itself.
The contribution is testing whether 483 text carries signal that the coarse
OAI/VAI/NAI outcome collapses away, particularly inside the VAI class, and
whether that signal predicts facility-level adverse events the way the
silent-problem hypothesis predicts (Section~\ref{sec:results}).
We validate the extraction pipeline through an external pharmaceutical-%
quality expert review and an independent blind-labeler round, and we build
a time-aware facility-level feature panel, cumulative text-signal snapshots
at each (FEI, inspection date) so no feature ever leaks future information,
plus a cross-inspection repeat detector that catches facilities cited for
the same deficiency years apart.
\note{Positioning against prior applied-LLM work: this sits in the tradition
of using a language model as a measurement instrument inside a larger
empirical design \citep{AppliedEconLLM2025, TAISR2024}, closest in spirit to
work that extracts structured text signal as input to a downstream outcome
model \citep{FusedLLMStartup2025, Li2026LLMDevice, AIatFDA2025}. Check:
author/year fields for these four are unverified placeholders in
references.bib, confirm before submission, and decide whether this
positioning belongs back in the visible text or stays here as a note.}

% ══════════════════════════════════════════════════════════════════════════
\section{Related Work}
\label{sec:lit}

\subsection{NLP and LLMs for FDA Regulatory Text}

Applying machine learning to FDA regulatory documents has concentrated on two
corpora: drug labels and adverse-event reports.
On drug labels, \citet{Bayer2021} evaluated 23 NLP systems for adverse drug
event extraction and found that the best systems reach F1 $\approx 0.89$,
suitable for human-assisted review but not fully autonomous processing.
\citet{Tatonetti2025OnSIDES} fine-tuned PubMedBERT on 200 curated labels to
produce 7.1 million drug--adverse-event pairs from 47,211 FDA and EMA labels
(F1 $= 0.90$; AUROC $= 0.92$).
\citet{Wu2021BERT} achieved Matthews Correlation Coefficients of 0.84--0.87
on drug-induced liver injury classification from FDA labels, compared to 0.60
for keyword rules.
A common finding across these studies is that fine-tuned domain-adapted models
(PubMedBERT, BioBERT) outperform general-purpose models on tasks requiring
specific biomedical vocabulary, but this advantage has narrowed sharply as
generative LLMs have grown in capability.
\citet{Li2026LLMDevice} demonstrated that GPT-4-class models extract structured
data from FDA medical device adverse-event reports (MDRs) and pre-market
summaries at 93--97\% agreement with human annotators, \emph{without any
fine-tuning on domain corpora}, a result that effectively subsumes the
PubMedBERT fine-tuning advantage for well-specified extraction tasks.
Similarly, \citet{Hassani2025} found GPT-4o achieves 89\% precision and 87\%
recall classifying regulatory provisions in food-safety regulations without
task-specific training, and \citet{AIatFDA2025} report comparable LLM
extraction performance directly on FDA regulatory and clinical-trial text.
\citet{Atalay2026RecallRisk} applied multi-task BERT to 54,165 FDA device
recall narratives (0.963 accuracy), confirming that FDA regulatory free text is
amenable to automated classification at scale.

On FDA inspection text specifically, \citet{Pazhayattil2020} conducted the
only published quantitative study of Form 483 data using pre-structured citation
frequency counts from a regulatory database (2014--2018, trend analysis of
observation rates by domain).
Warning-letter text has been analyzed descriptively for domain trends and data
integrity patterns \citep{Rathore2022, Park2025DataIntegrity, Kwiecinski2024},
but not used for feature extraction.
No prior study has applied NLP or LLM methods to the raw narrative text of Form
483 observations to extract structured quality-risk features.

We try to fill this gap.
Like \citet{Li2026LLMDevice}, we use a modern generative LLM with structured
JSON output instead of a fine-tuned BERT model.
Our signal-verification comparison against an independent classification
system (Redica, Section~\ref{sec:validation}) and our semantic-lift analysis
(Section~\ref{sec:results}) are both meant to show that a real share of
contamination and patient-risk signal needs contextual inference --
something keyword rules cannot do.

\subsection{Inspection Outcomes and Drug Quality Events}

The empirical relationship between FDA inspection activity and drug supply
disruptions is well-established.
\citet{Stomberg2018} provided the first regression-based evidence that
inspection and citation intensity are significantly associated with drug
shortage rates at the drug level.
\citet{Wang2025MSOM} extended this with instrumental variable analysis on US
manufacturing facilities: OAI inspection outcomes predict a 96.4\% lower
shortage likelihood in the subsequent 12 months, and VAI outcomes predict a
77.2\% reduction.
Interpreted alongside Stomberg, the pattern suggests that inspections trigger
corrective action that ultimately stabilizes supply, but the severity of the
underlying quality problem before inspection is not captured by binary OAI/VAI
classifications.

Machine learning models predict drug shortages from demand-side signals with
substantial accuracy.
\citet{Liu2021} achieved C-statistic $= 0.93$ on pharmacy observations using IV
status, generic-only market structure, and number of manufacturers.
\citet{Pall2023} achieved 69\% accuracy one month ahead using XGBoost on
pharmacy sales records for 784 Canadian drugs.
Both studies relied exclusively on market-structure and demand features; neither
incorporated manufacturing inspection or quality signals.
\citet{Kosmas2023} modeled FDA inspection timing as a POMDP to maximize
expected value of inspecting high-risk facilities; their model assumes quality
is revealed by inspection but cannot observe it from text.

Our paper adds the supply-side quality signal these models assume but
cannot observe: structured features pulled from the narrative content of
Form 483 observations.

\subsection{Adverse Events as a Manufacturing Quality Outcome}

FAERS adverse event reports are widely used for pharmacovigilance
disproportionality analysis at the drug level, but their potential as a
facility-level manufacturing quality metric has received limited attention.
\citet{Sardella2021} provided the conceptual foundation: pharmacovigilance is a
primary mechanism for detecting manufacturing-origin safety failures
post-approval, as demonstrated by the nitrosamine contamination crisis and the
heparin adulteration scandal, both of which were first surfaced through adverse
event signal clustering tracing to facility-level manufacturing decisions.
\citet{Brown2020} went further, calling for systematic use of FDA's Sentinel
and FAERS data to generate facility-level manufacturer quality signals, and
noting that the FDA already folds FAERS ``hazard signals'' into its risk-based
facility site-selection model for inspections. The agency itself treats
adverse event patterns as a quality-surveillance signal.

Using FAERS as a manufacturer-level quality comparator introduces potential
confounds.
\citet{Rahman2017FAERS} and \citet{Alatawi2018} demonstrated that FAERS
reporting rates differ between branded and generic products partly due to
perception bias (reporters over-attribute adverse events to generic products),
and developed an authorized-generic comparison design to control for this.
Our analysis largely sidesteps this confound because we compare facilities
manufacturing the \emph{same} API: inter-facility differences within a drug
control for drug-level perception bias.
Consistent with \citet{Potter2025}, we treat FAERS counts as a
\emph{predictive correlate} of manufacturing quality, not as a measure of
adverse event incidence or a causal indicator.

Critically, no prior study has used FEI-aggregated FAERS counts as the
dependent variable in a model of inspection-text quality signals, or validated
FAERS as a quality proxy against independent third-party chemical testing.
We contribute both: a cross-lag analysis of 483-text features against
facility-level FAERS counts at horizons of 0, 1, and 2 years post-inspection,
and external validation of the FAERS quality signal against Valisure
independent analytical testing results for the same 14 APIs (Spearman
$\rho = 0.86$, $p = 0.014$, $n = 7$ drugs).

\subsection{Positioning}

Table~\ref{tab:positioning} summarizes our contribution relative to the closest
prior work.

\begin{table}[h!]
\centering
\small
\caption{Positioning relative to prior work}
\label{tab:positioning}
\begin{tabular}{p{3.5cm} p{2.2cm} p{2.2cm} p{2.2cm} p{2.5cm}}
\toprule
& \textbf{Pazhayattil 2020} & \textbf{Wang et al.\ 2025} & \textbf{Liu/Pall 2021/23} & \textbf{This paper} \\
\midrule
483 observation text  & Counts only   & No  & No  & LLM extraction \\
Inspection outcome    & Citation freq & OAI/VAI  & No  & OAI/VAI + text \\
Quality outcome       & None          & Shortage & Shortage & FAERS + Valisure \\
Facility-level FEI    & Partial       & Yes      & No       & Yes \\
External validation   & No            & IV       & CV       & Valisure testing + blind human labeling \\
\bottomrule
\end{tabular}
\end{table}

% ══════════════════════════════════════════════════════════════════════════
\section{Data}
\label{sec:data}

We construct a unified regulatory event timeline for 129 facilities supplying
14 generic APIs, linked via FDA Facility Establishment Identifier (FEI) across
five sources yielding 1,280 total regulatory events (Table~\ref{tab:data}).

\begin{table}[h!]
\centering
\caption{Data sources and regulatory event coverage}
\label{tab:data}
\begin{tabular}{p{3.8cm} p{4.2cm} p{4.4cm}}
\toprule
\textbf{Source} & \textbf{Events / Coverage} & \textbf{Key Text Field} \\
\midrule
FDA Inspections + Citations & 1,019 inspections; 961 citations (127/129 FEIs) & CFR Short \& Long Description \\
Redica 483 Observations & 98 FEIs; 1,067 observations & Full observation narrative (pre-digitized) \\
FDA Warning Letters & 17 letters; 14 FEIs & Violation narrative; close-out status \\
FDA Drug Recalls & 49 API-matched; 18 FEIs & Reason for Recall (free text) \\
OASIS Import Refusals & 117 API-matched; 22 FEIs & Refusal charge codes; lab analysis flags \\
\bottomrule
\end{tabular}
\end{table}

Independent quality ground truth comes from Valisure analytical testing
scores for the same 14 APIs (2020--2024), used to externally validate the
FAERS quality signal (Section~\ref{sec:methods}).

\subsection{FDA Form 483 Observations}

FDA Form 483 is issued to facility management at the close of an inspection
when an investigator observes conditions that may violate the Food, Drug, and
Cosmetic Act (21 CFR Parts 210--211).
Each letter contains numbered observations written in investigator narrative
prose. They follow no templated structure beyond the opening CFR-paraphrase
sentence. That makes them the richest source of regulatory intelligence we
have, and also the least structured.

\paragraph{Why not FDA's public PDF portal directly.}
FDA only releases 483 letters as scanned PDF images, so any direct approach
needs OCR (e.g.\ Tesseract) before observation-level text is even available,
and OCR quality on scanned regulatory documents is uneven.
An early version of this project ran that pipeline and got usable text for
only 38 of 129 facilities, with error rates that hurt downstream extraction
on the documents with poor scan quality.
\note{Keep 1--2 sentences here on OCR accuracy/coverage as motivation --
this is background, not a data source we use going forward.}

\paragraph{Redica-provided text.}
We use pre-digitized 483 observation text licensed from Redica Systems
instead, a regulatory-intelligence vendor that keeps a structured database
of FDA inspection documents.
Redica's text extraction solves the OCR problem entirely: observation
narratives come as clean, machine-readable text, not scanned images, so we
get both better per-document text quality and much wider facility coverage
than the public-PDF route.
This gives full observation-level text for 98 of our 129 facilities: 1,067
individual observations across 246 inspections (2018--2026).
We use Redica's own pre-structured classifications (severity, violation
category, data-integrity flag) as a comparison point to check our LLM
extraction against an independent system, not as ground truth.
Redica's labels come from a different classification philosophy built for a
different purpose, so both agreement and disagreement tell us something
(Section~\ref{sec:methods}).

\subsection{Sample}

The 14 APIs are: Ampicillin, Ampicillin/Sulbactam, Atorvastatin, Bupropion,
Calcium Gluconate, Lisinopril, Magnesium Sulfate, Metformin, Metoprolol,
Metronidazole, Pantoprazole, Potassium Chloride, Tacrolimus, and Vancomycin.
FEI-to-drug mapping uses the Valisure API--FEI crosswalk (193 unique FEI--drug
pairs; 129 unique FEIs).

% ══════════════════════════════════════════════════════════════════════════
\section{Methodology}
\label{sec:methods}

We apply three layers of text analysis, progressing from rule-based NLP to
LLM extraction.

\paragraph{Layer 1: NLP on Structured Regulatory Text.}
CFR citation Short Descriptions (230 unique values across 961 citations) are
aggregated into six regulatory domains via CFR section mapping: Laboratory/QC
Controls, Production Controls, Sterile Manufacturing, Documentation/Data
Integrity, Facility/Equipment, and Materials/Components.
Recall reason free text is classified into nine categories (contamination,
cGMP failure, wrong potency, data integrity, etc.) using regex pattern matching.
Warning letter text is coded for violation domain and close-out resolution
status, providing a severity-escalation signal beyond routine citations.

\paragraph{Layer 2: LLM Pipeline for Form 483 Observation Extraction.}
Form 483 observations hold specific measurements, equipment identifiers,
and investigator narrative: the richest regulatory intelligence available,
but it comes as unstructured prose (Section~\ref{sec:data}).
The pipeline: Redica-provided observation text, then Claude (Anthropic)
extraction under a strict JSON schema, then field validation with
categorical coercion, structural self-consistency guards
(Section~\ref{sec:schema}), and an evidence guard.
Extraction is idempotent: an observation already scored is skipped on
re-run.
The prompt went through two calibration stages before any corpus-scale run.
First, iterative live testing on small samples (5--20 observations) to
settle on stable field definitions.
Second, a structured external review by a pharmaceutical-quality expert,
whose comments led to a revised prompt (v2): a restructured six-system
violation taxonomy, a split confirmed-vs.-plausible-risk contamination flag,
and a four-scenario patient-risk rule (Section~\ref{sec:schema}).
A blind human-labeler validation round against the v2 prompt is running now
(Section~\ref{sec:validation}).

\paragraph{Layer 3: Facility-Level Correlation Analysis.}
A per-facility feature matrix integrating all regulatory signals feeds
correlation analysis against Valisure quality scores and FAERS adverse
event outcomes.

\subsection{LLM Extraction Schema (Prompt v2)}
\label{sec:schema}

For each Redica-provided Form 483 observation we send the cleaned
observation text to a Claude model (Anthropic) with a structured JSON
schema prompt.
The prompt instructs the model to return the following fields:

\begin{itemize}[noitemsep]
  \item \textbf{Categorical}: violation category (six-system framework,
        following FDA's own Compliance Program Guidance Manual 7356.002
        drug-manufacturing inspection systems: QualitySystem,
        ProductionSystem, MaterialsSystem, FacilitiesEquipmentSystem,
        LaboratoryControlsSystem, PackagingLabelingSystem, or Other),
        severity tier (Critical/Major/Moderate/Minor), failure scope
        (SingleBatch/MultipleProducts/FacilityWide/Unclear), root-cause type
        (Capital/Cultural/Mixed/Unclear), remediation signal
        (Strong/Partial/Weak/None).
  \item \textbf{Binary flags}: repeat finding, patient risk, data integrity
        failure, investigation failure, and a \emph{mutually exclusive}
        confirmed-vs.-plausible contamination pair: \texttt{contamination\_flag}
        (an actual, confirmed contamination or sterility-failure event) and
        \texttt{contamination\_risk\_flag} (a contamination-control gap or
        risk with \emph{no} confirmed event). The two can never both be true
        on the same observation by construction; any residual concern about
        incomplete follow-up on a confirmed event is captured by the
        investigation-failure flag instead, not by re-using the same
        evidence for both contamination fields.
  \item \textbf{Qualitative}: severity rationale, root-cause rationale,
        patient-risk rationale, verbatim evidence quote, confidence score
        (0--1).
\end{itemize}

The severity tier is graded on \emph{documented product impact} rather than
the seriousness of the system failure: Critical requires that affected product
was released or distributed; Major requires an actual defect, failing result,
or unreliable data found on site; Moderate covers deficient procedures or
systems with no actual defect documented; Minor covers administrative gaps.
The patient-risk flag fires under one of four scenarios, calibrated to trade
off consequence severity against likelihood of patient exposure: (a) a
sterile/injectable dosage form with confirmed contamination. No release
evidence is required, because the injectable route bypasses the body's normal
filters and unit-level particulate detection is inherently incomplete;
(a2) a text-named high-risk drug class (narrow-therapeutic-index, oncology,
or nitrosamine) with a confirmed defect. Again, no release evidence is
required, because the consequence of even a contained deviation is large for
these classes; (b) a confirmed defect with stated release or distribution,
for any dosage form; or (c) product released without a documented QA
disposition, for any dosage form.
Scenarios (b) and (c) mean any dosage form, not only sterile/injectable
products, can trigger the flag; the asymmetry with (a)/(a2) is only in
how much textual evidence is required, not in which products are eligible.
The evidence quote is subject to an \emph{evidence guard}: the returned quote
must appear verbatim in the source text; otherwise it is cleared.

\subsection{Validation}
\label{sec:validation}

We calibrated the prompt in three stages, and results in this paper are
final only once all three are done.

\emph{Stage 1, live iterative calibration.}
We tested early drafts on small samples (5--20 observations) against four
automated checks: (a) the severity distribution is graded, not top-heavy;
(b) the Critical+Major share separates OAI from VAI inspections (via Redica
inspection classifications, matched on FEI and inspection date); (c) flag
fire rates land in an informative range; (d) the scope distribution spreads
across classes.
One early draft failed check (a): 84\% of sampled observations came back
Major. We revised it to the documented-product-impact rubric above.

\emph{Stage 2, external expert review.}
A pharmaceutical-quality expert reviewed the full prompt rule set and sent
back ten substantive comments. We addressed all of them in prompt v2:
restructured the violation-category taxonomy to FDA's own six-system CPGM
framework, split the contamination flag into a confirmed and a
plausible-risk version with a strict mutual-exclusivity rule, and added the
four-scenario patient-risk structure described above.

\emph{Stage 3, blind independent human-labeler validation (in progress).}
A trained research assistant, blind to model output, is labeling a
stratified pilot sample (n=50) drawn from the v2 extraction run.
Scoring against the model's labels is not done yet at the time of writing.
We plan a larger stratified round (target n$\approx$100--150, oversampling
rare positive classes) once we review the pilot.
\todo{Report Stage 3 agreement statistics (percent agreement, Cohen's
$\kappa$ per field) once available.}

We also compare LLM output against Redica's own independent pre-structured
classification of the same observation text, as a signal-verification
check, not a validation metric.
Redica is an independently built system with its own classification
philosophy, not a ground-truth label source, so we treat both agreement and
systematic disagreement as informative rather than as an accuracy score.
\todo{Report current signal-verification agreement statistics once the
comparison is re-run against the v2 prompt -- the existing comparison
predates the Stage 2 revision and is not valid for citation right now.}

\subsection{Facility-Level Aggregation}

We aggregate observation-level signals into \emph{time-aware cumulative
snapshots}: one row per (FEI, 483 inspection date), where each feature is
computed from all observations of that facility up to and including that
date.
The downstream AE model joins on FEI and takes the most recent snapshot
before each prediction period, so no feature ever encodes information from a
future inspection.
Rather than hand-weighted composite indices, raw feature shares are passed
directly to the model, which learns weights from data: severity tier shares
(including the combined Critical$+$Major share), scope shares, violation
category shares (six-system CPGM framework), root-cause and remediation
shares, the six LLM flag shares (repeat, patient risk, data integrity,
investigation failure, and the mutually-exclusive contamination-confirmed /
contamination-risk pair), and two regex flags with no LLM equivalent
(OOS/OOT events and prior-warning-letter references).

\paragraph{Cross-inspection repeat detection.}
The per-observation repeat flags only catch repeats the investigator
explicitly labels (``this is a repeat observation''), firing on 6.8\%
(regex) and 10.0\% (LLM) of observations.
We add a complementary signal that no per-observation method can produce:
each 483 observation opens with a standardized citation sentence (the
template text before ``Specifically,'').
When that sentence matches an observation from a strictly earlier inspection
of the same facility (token overlap $\geq 0.80$), the same deficiency was
cited again.
This flags 97 of 616 dated observations (15.7\%) as cross-inspection
repeats; the feature is time-safe by construction because it only looks
backward.
Full-text similarity fails at this task: shared GMP vocabulary gives
unrelated observations a $\sim$0.45 baseline token overlap, so matching is
restricted to the standardized citation sentence.

% Note to self: cut MQRI and shortage prediction from this paper's scope.
% Both may resurface as a separate paper later, but this draft stops at
% the AE analysis.

% ══════════════════════════════════════════════════════════════════════════
\section{Results}
\label{sec:results}

\begin{center}
\fbox{\parbox{0.92\textwidth}{
\textbf{PENDING, NOT CURRENT.} Every number in this section is from the
pre-Redica, pre-v2-prompt extraction run (622 observations, 38 facilities,
GPT-5-mini, old 8-class violation taxonomy, old contamination flag).
The Data and Methodology sections above now describe a different pipeline
(Redica-only corpus, 98 facilities, 1,067 observations, Claude v2 prompt,
six-system violation categories, mutually-exclusive contamination flags),
which has not yet been run to completion (Stage 3 blind validation with
Abdul is in progress; a full-corpus run is intentionally on hold until that
validation clears, to avoid paying for a run that a validation finding could
invalidate). \textbf{Nothing below this point should be cited.} It is kept
as a scaffold showing the intended shape of the analysis: table
structure, the two-failure-mode framing, the FAERS/Valisure validation
design. It is meant to be re-populated with current numbers once the pipeline output
is validated.
}}
\end{center}

\subsection{Descriptive Characterization of Form 483 Observations}

Table~\ref{tab:obs_desc} summarizes the LLM extraction results across
622 observations from 38 facilities.

\begin{table}[h!]
\centering
\caption{Observation-level signal distribution ($n = 622$ observations,
         38 facilities)}
\label{tab:obs_desc}
\begin{tabular}{lrr}
\toprule
\textbf{Signal} & \textbf{Count} & \textbf{Share (\%)} \\
\midrule
\multicolumn{3}{l}{\textit{Severity tier}} \\
\quad Critical  &  99 & 15.9 \\
\quad Major     & 188 & 30.2 \\
\quad Moderate  & 331 & 53.2 \\
\quad Minor     &   4 &  0.6 \\
\midrule
\multicolumn{3}{l}{\textit{Failure scope}} \\
\quad MultipleProducts & 374 & 60.1 \\
\quad FacilityWide     & 190 & 30.5 \\
\quad SingleBatch      &  46 &  7.4 \\
\quad Unclear          &  12 &  1.9 \\
\midrule
\multicolumn{3}{l}{\textit{Primary violation domain}} \\
\quad Production Controls  & 184 & 29.6 \\
\quad Lab Controls         & 159 & 25.6 \\
\quad Quality System       & 115 & 18.5 \\
\quad Buildings/Equipment  &  72 & 11.6 \\
\quad Records/Reports      &  53 &  8.5 \\
\quad Personnel            &  21 &  3.4 \\
\quad Packaging/Labeling   &  13 &  2.1 \\
\quad Other                &   5 &  0.8 \\
\midrule
\multicolumn{3}{l}{\textit{Root-cause type}} \\
\quad Cultural  & 271 & 43.6 \\
\quad Mixed     & 228 & 36.7 \\
\quad Capital   & 118 & 19.0 \\
\quad Unclear   &   5 &  0.8 \\
\midrule
\multicolumn{3}{l}{\textit{Remediation signal}} \\
\quad None    & 443 & 71.2 \\
\quad Partial & 103 & 16.6 \\
\quad Weak    &  74 & 11.9 \\
\quad Strong  &   2 &  0.3 \\
\bottomrule
\end{tabular}
\end{table}

\paragraph{Key finding 1: Severity grades on documented product impact.}
Under the documented-product-impact rubric, 15.9\% of observations are
Critical (defective product was released or distributed), 30.2\% Major
(an actual defect or failing result found on site), and 53.2\% Moderate
(deficient system, no defect documented).
The Critical share is striking: roughly one in six observations in this
sample documents defective product that reached the market.
The graded distribution gives the tier discriminating power, unlike a
rubric where any serious-sounding GMP failure rates highly:
the Critical$+$Major share is 20 percentage points higher in inspections
that FDA classified OAI than VAI.

\paragraph{Key finding 2: Cultural root causes dominate.}
80.3\% of observations are classified as Cultural or Mixed root cause,
suggesting that the predominant drivers of CGMP violations in this sample
are management, training, and quality-culture failures rather than
capital/equipment gaps.
This has direct implications for corrective action: capital investment alone
is insufficient.

\paragraph{Key finding 3: Remediation is rarely mentioned.}
71.2\% of observations mention no corrective action whatsoever, and only
two observations (0.3\%) reference a strong corrective action.
\note{This is partly structural: 483 observations are investigator-written
findings; firm responses are submitted separately.
However, Partial/Weak mentions do appear when investigators note commitments
made on-site. The absence is itself informative.}

\subsection{Semantic Lift: LLM vs.\ Regex}

Table~\ref{tab:lift} compares LLM flag rates to deterministic regex flags
on the same observations.

\begin{table}[h!]
\centering
\caption{Semantic lift: LLM vs.\ regex flag rates ($n = 622$ observations)}
\label{tab:lift}
\begin{tabular}{lrrr}
\toprule
\textbf{Flag} & \textbf{Regex share} & \textbf{LLM share}
              & \textbf{LLM-only lift} \\
\midrule
Contamination  & 33.1\% & 58.5\% & \textbf{27.3\%} \\
Patient risk   & 19.9\% & 33.6\% & \textbf{21.2\%} \\
Data integrity &  5.3\% & 18.5\% & 13.8\% \\
Investigation  & 37.5\% & 40.0\% &  7.2\% \\
Repeat finding &  6.8\% & 10.0\% &  3.2\% \\
\bottomrule
\end{tabular}
\end{table}

The contamination and patient-risk flags show the largest semantic lift:
the LLM flags 27 and 21 percentage points of observations, respectively,
that keyword rules miss.
This is because these inferences require contextual reasoning
(e.g., a sterile product exposed to unclassified air $\Rightarrow$ patient
risk, even though the words ``patient'' and ``contamination'' never appear)
that regex cannot replicate.
The reverse error also occurs: for patient risk, 7.6\% of observations are
regex-only: the word ``patient'' appears but the LLM judged that no
explicit harm pathway exists under the strict prompt definition.

\subsection{Facility-Level Text Features}

\todo{Table of top/bottom 10 FEIs by severity\_critmajor\_share and
repeat\_cross\_insp\_share with firm name (or anonymized).
Add: correlation of severity\_critmajor\_share with n\_oai.}

\subsection{Validation: FAERS as a Manufacturing Quality Proxy}

Before linking 483 text signals to adverse event outcomes, we validate that
FAERS serious AE counts track manufacturing quality rather than pharmacological
effects alone.
For each of the seven generic drugs with published Valisure analytical quality
scores, we compute the fraction of ANDA-holders with score $< 90$ (Valisure
failure rate) and the drug's FAERS serious adverse event count per million
Medicaid prescriptions (using CMS SDUD prescription volume as denominator).
Spearman correlation across seven drugs yields $\rho = 0.86$, $p = 0.014$:
drugs with higher proportions of analytically deficient products generate more
serious adverse events per patient exposed, even after controlling for
prescribing volume (Table~\ref{tab:valisure_val}).
The SDUD-normalized metric is stronger than the raw AE count ($\rho = 0.79$),
indicating that prescription volume was attenuating, not inflating, the
quality signal.

\begin{table}[h!]
\centering
\small
\caption{Valisure failure rate vs.\ FAERS serious AEs per million Medicaid Rx
         ($n = 7$ drugs; Spearman $\rho = 0.86$, $p = 0.014$)}
\label{tab:valisure_val}
\begin{tabular}{lrr}
\toprule
\textbf{Drug} & \textbf{Valisure fail rate} & \textbf{Serious AEs / M Rx} \\
\midrule
Tacrolimus         & 69\% & 1{,}738 \\
Metformin          & 48\% &   133 \\
Metoprolol         & 46\% &    61 \\
Potassium chloride & 41\% &    21 \\
Magnesium sulfate  &  0\% &    19 \\
Lisinopril         & 17\% &    16 \\
Calcium gluconate  & 33\% &     7 \\
\bottomrule
\end{tabular}
\end{table}

\subsection{Two Failure Modes: Governance vs.\ Technical Failures}

We test whether the same text signals that predict regulatory escalation also
predict patient harm, or whether inspection text encodes two distinct failure
pathways.
For each of 98 FEIs with 483 text features, we compute Spearman correlations
of each text signal against three outcomes: (1)~OAI inspection rate
(regulatory escalation), (2)~Class~I recall count (product-level regulatory
action), and (3)~mean FAERS serious AEs at lag $+$1~year (patient harm).

Results reveal a clear divergence (Table~\ref{tab:two_modes}):

\paragraph{Governance failures $\to$ regulatory action, not patient harm.}
\emph{Quality system violations} are the strongest predictor of OAI rate
across all ten text features ($\rho = +0.33$, $p < 0.001$) and show no
significant relationship with serious AEs ($\rho = -0.08$, $p > 0.4$).
Quality system observations cite inadequate oversight, missing review
procedures, and quality-unit authority failures, exactly the systemic
indicators FDA uses to justify OAI classification.
The regulatory response (mandatory CAPA, import alerts, product holds) limits
patient exposure, attenuating the harm signal.
\emph{Repeat observations} trend toward both OAI and recalls ($\rho \approx
+0.20$ for both), consistent with repeat citations building a regulatory
history that accelerates escalation.

\paragraph{Technical failures $\to$ patient harm, evading regulatory escalation.}
\emph{Lab controls violations} (inadequate test methods, OOS investigation
failures, missing specifications) are the strongest predictor of FAERS
serious AEs ($\rho = +0.32$, $p < 0.01$) and show essentially no relationship
with OAI rate ($\rho = +0.01$).
Lab controls are the gatekeeping function: when they fail, defective product
passes QC without generating the governance-level signals that trigger OAI.
\emph{Data integrity violations} follow the same pattern ($\rho = +0.22$ for
AEs, $p < 0.05$; $\rho = +0.10$ for OAI, $p > 0.3$).
DI failures are by design difficult to detect and slow to escalate.
Falsified records suppress the evidence trail that inspectors use to assign OAI.

\paragraph{Removed products.}
\emph{Contamination flags} show negative correlations with both OAI
($\rho = -0.05$) and AEs ($\rho = -0.19$), consistent with contamination
findings triggering recalls that remove product from the market before
widespread patient harm materializes.

\begin{table}[h!]
\centering
\small
\caption{Spearman $\rho$ — each 483 text signal vs.\ regulatory action and
         patient harm outcomes ($n = 98$ FEIs)}
\label{tab:two_modes}
\begin{tabular}{lrrr}
\toprule
\textbf{Text signal} & \textbf{OAI rate} & \textbf{Class~I recalls}
                     & \textbf{Serious AEs (lag+1)} \\
\midrule
\multicolumn{4}{l}{\textit{Governance failures → regulatory action}} \\
\quad Quality system violations & $+0.33$*** & $-0.02$ & $-0.08$ \\
\quad Repeat observations       & $+0.20$    & $+0.19$ & $+0.13$ \\
\quad Investigation failures    & $+0.17$    & $+0.04$ & $+0.13$ \\
\midrule
\multicolumn{4}{l}{\textit{Technical failures → patient harm}} \\
\quad Lab controls violations   & $+0.01$    & $+0.18$ & $+0.32$** \\
\quad Data integrity violations & $+0.10$    & $-0.09$ & $+0.22$* \\
\quad Scope: facility-wide      & $+0.10$    & $-0.10$ & $+0.10$ \\
\midrule
\multicolumn{4}{l}{\textit{Removed products}} \\
\quad Contamination flag        & $-0.05$    & $-0.07$ & $-0.19$ \\
\quad Patient risk flag         & $+0.14$    & $-0.03$ & $-0.06$ \\
\bottomrule
\multicolumn{4}{l}{\footnotesize * $p<0.05$; ** $p<0.01$; *** $p<0.001$} \\
\end{tabular}
\end{table}

\subsection{Predictive Validity: Text Features vs.\ Valisure Outcomes}

\todo{Run correlation: severity\_critmajor\_share vs.\ Valisure dissolution
difference factor (2024).
Run correlation: contamination\_llm\_share vs.\ Valisure NDMA/DMF.}

% Note to self: cut the shortage-prediction benchmark subsection, table, and
% feature-importance figure. Out of scope for this paper -- stops at AE.

% ══════════════════════════════════════════════════════════════════════════
\section{Discussion}
\label{sec:discussion}

\begin{center}
\fbox{\parbox{0.92\textwidth}{
\textbf{PENDING, NOT CURRENT.} This section interprets the stale Results
above and inherits the same caveat: every specific number and claim here is
from the pre-Redica, pre-v2-prompt run and must be re-derived once the
current pipeline's output is validated. Kept as a scaffold for argument
structure, not as citable interpretation.
}}
\end{center}

\paragraph{Form 483 text as a leading quality signal.}
LLM classification of 483 observation narratives produces signals that
structured compliance databases do not capture: graded severity that tracks
FDA inspection outcomes (Critical$+$Major share 20pp higher in OAI than VAI
inspections), failure scope, remediation commitment, and chronic-violator
identification via cross-inspection repeat detection (15.7\% of
observations repeat citation language from an earlier inspection of the
same facility).
The semantic lift analysis shows that LLMs detect contamination-control risk
and patient-risk implications in 27 and 21 percentage points more
observations, respectively, than keyword rules. Contextual reasoning, not
pattern matching, is required to extract these signals.

\paragraph{Root cause composition.}
\note{Dropped the Agency Theory / Remediation Economics framing here (and
from Abstract/Intro) -- root\_cause\_type and remediation\_signal stay as
features, but we are not confident enough in prior runs' predictive power
for these fields to hang a theory on them. Revisit this paragraph once
current-pipeline numbers exist; may cut entirely or reframe as a purely
descriptive finding.}
80.3\% of observations carry Cultural or Mixed root causes under the old
pipeline (number to re-check). If cultural-root-cause facilities do show
different recovery dynamics than capital-root-cause ones, that is worth
reporting descriptively, without claiming it as evidence for a specific
theory.
\todo{Add survival analysis results when available.}

\paragraph{Open questions for the full analysis.}
Two questions structure where this line of work goes next.
First, does the silent-problem pattern hold up across the full corpus and
across drugs, not just in the descriptive cases shown here?
Second, can we say anything about why some facilities with alarming text end
up VAI while others end up OAI, given that the EIR and the facility's
response are not observable to us?
Shortage prediction and a composite quality risk index are natural
extensions of this work, but sit outside the scope of the current paper.

\paragraph{Limitations.}
\begin{itemize}[noitemsep]
  \item \textit{Coverage}: Redica text is available for 98 of 129 FEIs;
        structured citation signals cover 127/129.
        Drug-level text indices reflect only the subset of FEIs with
        available inspection text, so a drug where most FEIs lack text
        coverage will have a thinner, less representative average.
  \item \textit{Cross-inspection repeat detection}: the template-sentence
        matcher catches deficiencies re-cited under the same standardized
        citation language (e.g., FEI 3011905047, a sterile injectable
        manufacturer, was cited for the identical employee-training
        deficiency in 2019 and again in 2024).
        It will miss repeats that an investigator cites under a different
        CFR paraphrase, so the 15.7\% cross-inspection repeat rate is a
        lower bound.
  \item \textit{Sample size}: 14 drugs and a modest count of OAI inspections
        yield wide confidence intervals; all findings are exploratory.
  \item \textit{LLM reliability}: manual review of a sample confirms strong
        precision on violation category and severity; contamination flags
        occasionally over-trigger on training-context observations.
  \item \textit{Selection}: Valisure tests a non-random subset of generics.
  \item \textit{Causality}: we estimate predictive associations, not causal effects.
  \item \textit{The EIR gap}: outcome assignment depends on facility
        responses we cannot observe, so we cannot fully explain why a given
        facility received the classification it did.
\end{itemize}

% ══════════════════════════════════════════════════════════════════════════
\section{Conclusion}
\label{sec:conclusion}

\begin{center}
\fbox{\parbox{0.92\textwidth}{
\textbf{PENDING, NOT CURRENT.} The specific findings restated below are
from the stale pre-Redica, pre-v2-prompt run (Section~\ref{sec:results}
caveat applies identically here) and must be replaced once the current
pipeline output is validated. The framing paragraph and applications
paragraph are less number-dependent and closer to final.
}}
\end{center}

We present a regulatory intelligence framework that integrates Redica's
483 observation text, structured FDA databases, and an LLM extraction
pipeline to turn inspector narrative into structured, patient-harm-relevant
signals.
Key findings, pending the current-pipeline re-run (Section~\ref{sec:results}
caveat applies): (1) LLM-graded severity tracks real inspection outcomes,
and a meaningful share of observations document defective product that
reached the market; (2) cross-inspection repeat detection identifies
chronic violators that per-observation flags miss; (3) contamination and
patient-risk signals require contextual reasoning that keyword rules
cannot replicate.
The central open question is the silent-problem pattern: whether 483 text
predicts elevated adverse events specifically among VAI-classified
facilities, where the outcome label alone cannot tell a genuinely resolved
problem apart from an unresolved one.
The pipeline is automated and reproducible from licensed and publicly
available FDA data.
If the silent-problem pattern holds, it points to a scalable use case for
manufacturing quality surveillance: catching the facilities where
enforcement action stopped short of what the underlying problem warranted.
To our knowledge, this is the first study to link FDA regulatory inspection
text directly to facility-level adverse event outcomes.

% ══════════════════════════════════════════════════════════════════════════
\bibliographystyle{agsm}
\bibliography{references}

% ══════════════════════════════════════════════════════════════════════════
\appendix

\section{LLM Flag Definitions (Prompt v2)}
\label{app:flags}

\paragraph{Extraction model.} Claude (Anthropic) with structured JSON output
schema (strict mode); model selection informed by a three-way comparison
against Claude Haiku 4.5, Claude Sonnet 5, and GPT-5-mini on a common
50-observation sample. \todo{State final production model once Stage 3
validation (Section~\ref{sec:validation}) is complete.}
Extraction is idempotent: failed or already-scored rows are re-scored on the
next run rather than silently imputed.

\paragraph{Evidence guard.} The \texttt{evidence\_quote} field must appear
verbatim in the source observation text. If the model paraphrases, the field
is cleared to empty string rather than accepting an inaccurate citation.

\paragraph{Violation category (six-system framework).} Following FDA's
Compliance Program Guidance Manual 7356.002 (drug manufacturing inspections),
not the medical-device QSIT framework used in an earlier prompt draft:
QualitySystem, ProductionSystem, MaterialsSystem, FacilitiesEquipmentSystem,
LaboratoryControlsSystem, PackagingLabelingSystem, or Other.

\paragraph{Flag definitions} are as specified in the extraction prompt
(see \texttt{01\_extract\_observation\_signals.py}, \texttt{--prompt-version
v2}):

\begin{description}[leftmargin=2em]
  \item[\texttt{repeat\_flag\_llm}] True only when the observation explicitly
    states this is a repeat or previously-cited finding from a prior inspection.
    Not true for repeated examples within the same current observation.
  \item[\texttt{patient\_risk\_flag\_llm}] True only when one of four explicit
    harm-pathway scenarios applies: (a) sterile/injectable product with
    confirmed contamination, (a2) a text-named high-risk drug class (NTI,
    oncology, nitrosamine) with a confirmed defect, (b) a confirmed defect
    with stated release or distribution (any dosage form), or (c) product
    released without a documented QA disposition (any dosage form).
    Scenarios (a)/(a2) require no release evidence; (b)/(c) do.
    ``Could affect quality'' does not qualify under any scenario.
  \item[\texttt{data\_integrity\_flag\_llm}] True only for explicit data
    trustworthiness failures (falsification, backdating, audit-trail problems).
    Not true for ordinary incomplete documentation.
  \item[\texttt{contamination\_flag\_llm}] True only for an ACTUAL, CONFIRMED
    contamination or sterility/cross-contamination event (microbial growth
    found, particulate matter observed in product, a confirmed sterility
    test failure, or confirmed cross-contamination in product).
  \item[\texttt{contamination\_risk\_flag\_llm}] True only when NO
    contamination is confirmed anywhere in the observation, and a
    contamination-control gap is described instead (sterility assurance
    failures, aseptic processing deficiencies, EM excursions without confirmed
    contamination, inadequate cleaning/sanitization, cross-contamination
    controls; list illustrative, not exhaustive). \textbf{Mutually exclusive
    with \texttt{contamination\_flag\_llm} by construction}: false
    whenever that flag is true, with no exception, even if the follow-up
    investigation on a confirmed event was incomplete (that belongs to
    \texttt{investigation\_flag\_llm} instead, not this flag; the two fields
    must not be justified by the same evidence).
  \item[\texttt{investigation\_flag\_llm}] True only for an explicit failed,
    missing, delayed, or inadequate investigation of a concrete event
    (deviation, OOS/OOT, contamination event, batch failure). Not true for a
    generic missing-procedure statement with no specific event attached.
\end{description}

\paragraph{Removed in prompt v2.} The former \texttt{systemic\_flag\_llm}
fired on 94\% of observations (uninformative) and was replaced by the
4-class \texttt{scope} field; \texttt{documentation\_flag\_llm} is covered
by the RecordsReports-equivalent violation category.

\paragraph{Changed in prompt v2 (external expert review).} Violation
category restructured from an 8-class ad hoc taxonomy to FDA's own
six-system CPGM 7356.002 framework; the single \texttt{contamination\_flag}
split into the confirmed/plausible-risk pair above; \texttt{patient\_risk\_flag}
expanded from a 3-scenario to a 4-scenario rule (adding the named-high-risk-%
drug-class scenario).

\section{Severity Tier Definitions (Prompt v2)}
\label{app:severity}

The tier is decided by one question: what level of actual product impact
does the text document?

\begin{description}[leftmargin=2em]
  \item[Critical] Affected product was released or distributed (e.g.,
    contaminated lots distributed before the investigation was closed).
  \item[Major] An actual defect, failure, or unreliable result found at the
    facility but no evidence of release (e.g., particulate matter observed
    in lots; test results invalidated without quality approval).
  \item[Moderate] A deficient procedure, system, or practice with no actual
    product defect documented (e.g., media fills that do not mimic routine
    production; systems that allow data deletion without evidence it
    occurred). The default tier for most observations.
  \item[Minor] Documentation or administrative gap with no plausible product
    impact.
\end{description}

\end{document}
